\documentclass[11pt]{article}

\usepackage[margin=1.1in]{geometry}
\usepackage{amsmath, amssymb}
\usepackage{booktabs}
\usepackage{array}
\usepackage{xcolor}
\usepackage{listings}
\usepackage{hyperref}
\usepackage{microtype}
\usepackage{parskip}
\usepackage{titlesec}
\usepackage{enumitem}
\usepackage{fancyhdr}
\usepackage{caption}
\usepackage{subcaption}
\usepackage{float}
\usepackage{mathtools}

\hypersetup{
  colorlinks=true,
  linkcolor=black,
  urlcolor=blue!60!black,
  citecolor=black
}

\lstset{
  basicstyle=\ttfamily\small,
  backgroundcolor=\color{gray!8},
  frame=single,
  framerule=0.4pt,
  rulecolor=\color{gray!40},
  breaklines=true,
  captionpos=b,
  language=Python,
  keywordstyle=\bfseries\color{blue!70!black},
  commentstyle=\itshape\color{gray!60!black},
  stringstyle=\color{green!50!black},
  numbers=left,
  numberstyle=\tiny\color{gray},
  numbersep=6pt,
}

\pagestyle{fancy}
\fancyhf{}
\rhead{\small NFL Analytics Framework}
\lhead{\small Vasugov}
\cfoot{\thepage}
\renewcommand{\headrulewidth}{0.4pt}

\titleformat{\section}{\large\bfseries}{{\thesection.}}{0.5em}{}
\titleformat{\subsection}{\normalsize\bfseries}{{\thesubsection}}{0.5em}{}
\titleformat{\subsubsection}{\normalsize\itshape}{{\thesubsubsection}}{0.5em}{}

\newcommand{\E}{\mathbb{E}}
\newcommand{\R}{\mathbb{R}}

\begin{document}

% ─────────────────────────────────────────────────────────────────────────────
%  TITLE
% ─────────────────────────────────────────────────────────────────────────────

\begin{titlepage}
\centering
\vspace*{2cm}

{\LARGE\bfseries NFL Analytics Framework}\\[0.6em]
{\large A Gradient Boosting Approach to\\Play-Level Expected Points Modelling}\\[2.5cm]

{\large Vasu Gov}\\[0.4em]
{\normalsize \texttt{vasugov/analytics}}\\[2cm]

{\normalsize 2025--2026}\\[1cm]

\vfill
\begin{abstract}
\noindent
This paper documents the design, implementation, and results of an NFL analytics framework
that moves beyond descriptive aggregate statistics toward predictive, model-derived metrics.
An R ingestion pipeline loads play-by-play data from nflfastR across the 2021--2023 NFL
regular seasons, producing 101,474 labelled pass and run plays. Python feature engineering
transforms raw game-state variables into 27 model inputs including interaction terms,
situational flags, and normalised field-position features. Five XGBoost gradient boosting
models are trained using a strictly time-ordered cross-validation scheme (2021--22 training,
2023 holdout): an EPA regression model, a play-success binary classifier, a red-zone
touchdown probability model, a win probability model, and a drive-outcome multi-class
classifier. Trained models are serialised and their outputs are exported as JSON to a
static GitHub Pages web application offering an interactive league-wide ranking table,
team profile charts, feature importance visualisations, and a real-time play predictor.
The paper explains every design decision in depth---from the mathematics of gradient
boosting to the rationale for time-based rather than random cross-validation---and
contextualises the model results within the inherent noise floor of play-level prediction.
\end{abstract}
\end{titlepage}

\tableofcontents
\newpage

% ─────────────────────────────────────────────────────────────────────────────
%  1. INTRODUCTION
% ─────────────────────────────────────────────────────────────────────────────
\section{Introduction}

The modern NFL analytics revolution began in earnest with the public release of
play-by-play data and, later, the \texttt{nflfastR} R package, which provides
pre-computed Expected Points (EP) and Win Probability (WP) estimates for every
play in the NFL since 1999. These pre-built metrics democratised sophisticated
football analysis but left a gap: they treat team performance at the aggregate
level and offer limited tools for building \emph{custom predictive models} that
can answer novel questions.

This project fills that gap. It constructs an end-to-end machine learning pipeline
that:
\begin{enumerate}[noitemsep]
  \item ingests raw play-by-play data from nflfastR,
  \item engineers a rich play-level feature dataset,
  \item trains five XGBoost gradient boosting models on those features,
  \item serves predictions via a FastAPI backend, and
  \item publishes results to a static web dashboard.
\end{enumerate}

The central modelling target is \emph{Expected Points Added} (EPA)---the most
important single-play metric in modern football analytics. By learning the mapping
from game state to EPA, we can ask: given the down, distance, field position, score,
and time remaining, what is the \emph{expected} change in scoring outcome this play
will produce? Teams whose actual EPA consistently exceeds model expectation are
demonstrably outperforming their situational context.

The remainder of this paper is structured as follows. Section~\ref{sec:background}
introduces the football metrics used throughout. Section~\ref{sec:data} describes the
dataset. Section~\ref{sec:features} documents feature engineering. Section~\ref{sec:gb}
explains gradient boosting theory. Section~\ref{sec:models} specifies each model.
Section~\ref{sec:training} covers training methodology. Section~\ref{sec:results}
presents empirical results. Section~\ref{sec:system} describes the full system
architecture. Section~\ref{sec:webapp} documents the web application. Sections
\ref{sec:discussion} and~\ref{sec:future} discuss limitations and future directions.

% ─────────────────────────────────────────────────────────────────────────────
%  2. BACKGROUND
% ─────────────────────────────────────────────────────────────────────────────
\section{Background}
\label{sec:background}

\subsection{Expected Points (EP)}

The concept of Expected Points formalises the intuition that field position has
monetary value in terms of future scoring. Carter and Machol (1971) first computed
EP empirically from NFL data. The idea is straightforward: for any game state
$(d, y, l)$---down $d$, yards to go $y$, and yard line $l$ (yards to the end zone)
---define EP as the expected number of points the possession team will score on
this drive minus the points they will concede before regaining possession:
\begin{equation}
  \text{EP}(d, y, l) = \E[\text{next scoring event} \mid d, y, l]
\end{equation}
where each scoring event is signed relative to the possession team: a touchdown
scores $+7$, a safety concedes $-2$, etc. Modern implementations (nflfastR) extend
this formulation to include score differential, time remaining, and timeout counts
as additional state variables.

\subsection{Expected Points Added (EPA)}

EPA is the change in EP produced by a single play:
\begin{equation}
  \text{EPA}_t = \text{EP}(s_{t+1}) - \text{EP}(s_t)
\end{equation}
where $s_t$ is the full game state immediately before play $t$ and $s_{t+1}$ is the
state immediately after. A ten-yard gain that converts a third down produces positive
EPA; an incomplete pass that wastes a down produces negative EPA. The key property is
that EPA is \emph{context-adjusted}---a five-yard gain on 3rd \& 4 is worth much more
than the same gain on 1st \& 10.

\subsection{Play Success}

A play is defined as \emph{successful} if and only if $\text{EPA} > 0$:
\begin{equation}
  \text{success}_t = \mathbf{1}[\text{EPA}_t > 0]
\end{equation}
Success rate---the fraction of plays with positive EPA---is a stable, easily
interpretable team-level offensive metric.

\subsection{Win Probability Added (WPA)}

The Win Probability function $\text{WP}(s)$ maps the current game state to the
probability that the possession team wins the game. WPA is the play-level change:
\begin{equation}
  \text{WPA}_t = \text{WP}(s_{t+1}) - \text{WP}(s_t)
\end{equation}
Cumulative WPA summarises how much each team contributed to (or detracted from)
their win chances over the course of a season.

\subsection{Red Zone Efficiency}

``Red zone'' refers to plays within the opponent's 20-yard line ($l \leq 20$).
Red zone efficiency is typically measured as the fraction of red-zone possessions
that result in touchdowns. This is a high-leverage metric: the difference between
scoring a touchdown and a field goal (or nothing) is three to seven points.

\subsection{Drive Efficiency}

A drive is the sequence of consecutive plays by one team. Drive efficiency captures
the outcome distribution at the drive level---what fraction of drives end in
touchdowns, field goals, punts, or turnovers---aggregating away play-level noise.

% ─────────────────────────────────────────────────────────────────────────────
%  3. DATA
% ─────────────────────────────────────────────────────────────────────────────
\section{Data}
\label{sec:data}

\subsection{nflfastR Play-by-Play}

The nflfastR R package provides play-by-play data for every NFL game since 1999,
including pre-computed EP and WP values derived from a large multinomial logistic
regression model trained on the full historical record. Each row represents one
play and includes over 350 columns covering game identifiers, player names, game
state, pre- and post-play probabilities, and outcomes.

This project uses three regular seasons: 2021, 2022, and 2023. Regular-season
filtering is applied via \texttt{season\_type == "REG"}, excluding pre-season and
playoff games which have different strategic dynamics.

\subsection{Aggregate Metric CSVs}

The R pipeline (\texttt{R/pipeline/run\_pipeline.R}) computes five team-season-level
summary tables and writes them to \texttt{data/outputs/}:

\begin{table}[H]
\centering
\caption{Aggregate output files}
\begin{tabular}{llr}
\toprule
File & Columns & Rows \\
\midrule
\texttt{epa\_by\_team.csv} & season, posteam, plays, total\_epa, epa\_per\_play & 97 \\
\texttt{success\_rate\_by\_team.csv} & season, posteam, plays, success\_rate & 97 \\
\texttt{wpa\_by\_team.csv} & season, posteam, total\_wpa, avg\_wpa & 100 \\
\texttt{redzone\_efficiency.csv} & season, posteam, rz\_plays, rz\_tds, rz\_td\_rate & 97 \\
\texttt{drive\_efficiency.csv} & season, posteam, drives, td\_drives, scoring\_rate, avg\_drive\_epa & 100 \\
\bottomrule
\end{tabular}
\end{table}

\subsection{Play-Level Feature Dataset}

For model training, the pipeline also exports \texttt{play\_features.csv} via
\texttt{R/metrics/features.R}. This file contains one row per pass or run play,
filtered to rows where down, yards to go, yard line, EPA, and time are all
non-missing. The dataset contains \textbf{101,474 plays} across three seasons.

Filtering criteria applied in R:
\begin{itemize}[noitemsep]
  \item \texttt{play\_type \%in\% c("pass", "run")} --- exclude special teams, kneel-downs
  \item Non-NA down, ydstogo, yardline\_100, epa, half\_seconds\_remaining, game\_seconds\_remaining
\end{itemize}

Season breakdown:
\begin{table}[H]
\centering
\caption{Play counts by season}
\begin{tabular}{lr}
\toprule
Season & Plays \\
\midrule
2021 & 33,412 \\
2022 & 33,981 \\
2023 & 34,081 \\
\midrule
Total & 101,474 \\
\bottomrule
\end{tabular}
\end{table}

% ─────────────────────────────────────────────────────────────────────────────
%  4. FEATURE ENGINEERING
% ─────────────────────────────────────────────────────────────────────────────
\section{Feature Engineering}
\label{sec:features}

All feature transformations live in \texttt{python/features/engineering.py}.
The \texttt{add\_features(df)} function takes a raw play DataFrame and returns
an enriched copy. The function is pure (no in-place mutation), making it safe
to call in notebooks, the trainer, and the API.

\subsection{Base Game-State Features}

These are the raw columns from the play-level CSV that directly encode the game
state at the moment of each snap:

\begin{table}[H]
\centering
\caption{Base game-state features}
\begin{tabular}{lp{9cm}}
\toprule
Feature & Description \\
\midrule
\texttt{down} & Current down (1--4) \\
\texttt{ydstogo} & Yards needed for a first down \\
\texttt{yardline\_100} & Yards to the opponent's end zone (1--99) \\
\texttt{score\_differential} & posteam score minus defteam score \\
\texttt{half\_seconds\_remaining} & Seconds left in the current half \\
\texttt{game\_seconds\_remaining} & Seconds left in the game \\
\texttt{posteam\_timeouts\_remaining} & Offensive team's timeouts (0--3) \\
\texttt{defteam\_timeouts\_remaining} & Defensive team's timeouts (0--3) \\
\texttt{is\_two\_minute} & 1 if within two-minute warning of a half \\
\texttt{is\_final\_two} & 1 if within final two minutes of the game \\
\texttt{is\_4th\_quarter} & 1 if in the fourth quarter \\
\texttt{is\_red\_zone} & 1 if yardline\_100 $\leq 20$ \\
\texttt{is\_goal\_to\_go} & 1 if yards to go equals yards to end zone \\
\texttt{shotgun\_flag} & 1 if play was from shotgun formation \\
\texttt{no\_huddle\_flag} & 1 if no-huddle pace was used \\
\texttt{play\_type\_bin} & 1 for pass, 0 for run \\
\bottomrule
\end{tabular}
\end{table}

\subsection{Engineered Features}

Beyond the raw state, several derived features capture relationships that
tree models might otherwise require many splits to approximate:

\begin{table}[H]
\centering
\caption{Engineered features}
\begin{tabular}{lp{9cm}}
\toprule
Feature & Rationale \\
\midrule
\texttt{down\_x\_ydstogo} $= d \times y$ & Down and distance interact; 3rd \& 8 is qualitatively different from 1st \& 8 \\
\texttt{yardline\_x\_down} $= l \times d$ & Field position matters more in late downs \\
\texttt{score\_x\_time} $= s \cdot (g/3600)$ & Large deficits matter more early; zero late \\
\texttt{ydstogo\_clipped} $= \min(y, 20)$ & Diminishing marginal difficulty beyond 20 yards \\
\texttt{field\_position\_norm} $= (100 - l)/100$ & Normalised progress toward end zone $\in [0,1]$ \\
\texttt{score\_abs} $= |s|$ & Magnitude of score imbalance \\
\texttt{time\_pressure} $= |s| / (g/3600 + 0.01)$ & Urgency score: large deficit, little time \\
\texttt{ydstogo\_bucket} & \texttt{np.digitize} bin of yards to go: short / medium / standard / long / very long \\
\texttt{field\_zone} & \texttt{np.digitize} bin of field position: goal line through deep own territory \\
\texttt{score\_bucket} & \texttt{np.digitize} bin of score differential \\
\texttt{quarter\_bucket} & \texttt{np.digitize} bin of game time \\
\texttt{score\_x\_time} & $s \cdot (g / 3600)$ \\
\bottomrule
\end{tabular}
\end{table}

\textbf{Implementation note on binning.} An early version of this code used
\texttt{pandas.cut()} with integer labels to create bin features.
\texttt{pd.cut} with labels returns a \texttt{Categorical} dtype column.
XGBoost 3.x uses the DataFrame Interchange Protocol (\texttt{\_\_dataframe\_\_})
when consuming a \texttt{pandas.DataFrame}; this protocol exposes
\texttt{Categorical} columns with a 2D shape representation that triggers
a shape-check assertion in \texttt{xgboost/data/array\_interface.h}.
The fix is to use \texttt{numpy.digitize()}, which returns a plain
\texttt{int64} array, and to convert all DataFrames to numpy arrays
(via \texttt{.to\_numpy()}) before passing them to any XGBoost method.

\subsection{Temporal Train/Validation Split}

The \texttt{time\_split(df, val\_season=2023)} function partitions the dataset:
\begin{align*}
  \mathcal{D}_{\text{train}} &= \{(x_t, y_t) : \text{season}(t) < 2023\} \quad (67,393 \text{ plays}) \\
  \mathcal{D}_{\text{val}} &= \{(x_t, y_t) : \text{season}(t) = 2023\} \quad (34,081 \text{ plays})
\end{align*}
The rationale for this chronological split is discussed in
Section~\ref{sec:training}.

% ─────────────────────────────────────────────────────────────────────────────
%  5. GRADIENT BOOSTING
% ─────────────────────────────────────────────────────────────────────────────
\section{Gradient Boosting Theory}
\label{sec:gb}

\subsection{Boosting: Ensembles of Weak Learners}

Boosting is an ensemble method that constructs a strong predictor by
sequentially adding weak predictors, each correcting the errors of the
previous ensemble. Given a loss function $\mathcal{L}$ and a set of base
learners $\{h_m\}$, the ensemble prediction at iteration $M$ is:
\begin{equation}
  \hat{F}_M(x) = \sum_{m=0}^{M} \eta \cdot h_m(x)
\end{equation}
where $\eta$ is the learning rate (shrinkage parameter). At each step $m$, we
fit $h_m$ to the \emph{negative gradient} of the loss evaluated at the current
prediction:
\begin{equation}
  r_i^{(m)} = -\left[ \frac{\partial \mathcal{L}(y_i, F(x_i))}{\partial F(x_i)} \right]_{F = \hat{F}_{m-1}}
\end{equation}
This is gradient descent in \emph{function space}: rather than updating parameters,
we update the entire prediction function in the direction that most reduces the loss.

\subsection{XGBoost: Regularised Gradient Boosting Trees}

XGBoost (Chen \& Guestrin, 2016) extends gradient boosting in two important
ways: (1) it uses a second-order Taylor expansion of the loss for a more
accurate descent direction, and (2) it adds explicit regularisation to the
objective.

The XGBoost objective at iteration $m$ is:
\begin{equation}
  \mathcal{L}^{(m)} = \sum_{i=1}^{n} \left[
    g_i h_m(x_i) + \frac{1}{2} h_i^{(m)} h_m(x_i)^2
  \right] + \Omega(h_m)
\end{equation}
where $g_i = \partial_{\hat{y}} \mathcal{L}(y_i, \hat{y})$ and
$h_i = \partial_{\hat{y}}^2 \mathcal{L}(y_i, \hat{y})$ are the first and second
derivatives of the loss with respect to the current prediction, and:
\begin{equation}
  \Omega(h) = \gamma T + \frac{1}{2} \lambda \sum_{j=1}^{T} w_j^2 + \alpha \sum_{j=1}^T |w_j|
\end{equation}
penalises tree complexity: $T$ is the number of leaves, $w_j$ are leaf weights,
$\gamma$ is the minimum gain threshold for splitting, $\lambda$ is the L2
regularisation parameter, and $\alpha$ is the L1 parameter.

The optimal leaf weight for leaf $j$ given the second-order approximation is:
\begin{equation}
  w_j^* = -\frac{\sum_{i \in \mathcal{I}_j} g_i}{\sum_{i \in \mathcal{I}_j} h_i + \lambda}
\end{equation}
where $\mathcal{I}_j$ is the set of training examples in leaf $j$.

\subsection{Why XGBoost for NFL Data?}

XGBoost (and gradient boosting in general) is well-suited to this problem for
several reasons:
\begin{itemize}
  \item \textbf{Tabular data.} NFL plays are naturally tabular: a fixed set of
    numeric and binary features with no spatial or temporal sequence structure
    within a single play. For such data, tree ensembles consistently outperform
    neural networks at moderate dataset sizes.
  \item \textbf{Feature interactions.} Football strategy is deeply contextual.
    Down and distance interact; score and time interact. Trees discover these
    interactions automatically through their recursive splitting structure.
  \item \textbf{Missing data.} Some game-state fields are occasionally absent.
    XGBoost handles missing values natively by learning which branch is better
    for missing observations during tree construction.
  \item \textbf{Interpretability.} XGBoost provides feature importance scores
    (gain, cover, frequency) and integrates with SHAP for local explanations,
    which is essential for understanding \emph{why} the model predicts what it does.
  \item \textbf{Efficiency.} The histogram-based tree construction (\texttt{tree\_method="hist"})
    bins continuous features, dramatically reducing split-finding cost. Training
    500--600 trees on 70,000 plays completes in under five seconds on a laptop.
\end{itemize}

% ─────────────────────────────────────────────────────────────────────────────
%  6. MODELS
% ─────────────────────────────────────────────────────────────────────────────
\section{Model Specifications}
\label{sec:models}

All models inherit from the \texttt{NFLModel} base class
(\texttt{python/models/base.py}), which provides a unified interface:
\texttt{fit\_with\_eval()}, \texttt{predict()}, \texttt{predict\_proba()},
\texttt{evaluate()}, \texttt{save()}, and \texttt{load()}.

\subsection{EPA Model (Regression)}

\textbf{Target:} $\text{EPA}_t \in \R$ (continuous, pre-computed by nflfastR).\\
\textbf{Loss:} Mean squared error $\mathcal{L} = (y - \hat{y})^2$.\\
\textbf{Features:} All 27 game-state and engineered features.

The EPA model learns the conditional expectation:
\begin{equation}
  \hat{\text{EPA}}(x) = \E[\text{EPA} \mid \text{game state} = x]
\end{equation}
This is the most fundamental model in the framework. Its residuals
$\epsilon_i = \text{EPA}_i - \hat{\text{EPA}}(x_i)$
form the basis for measuring \emph{offensive alpha}: how much a team's
actual EPA exceeds what the model predicts for their average game state.

\textbf{Hyperparameters:}
\begin{table}[H]
\centering
\begin{tabular}{ll}
\toprule
Parameter & Value \\
\midrule
\texttt{n\_estimators} & 600 \\
\texttt{max\_depth} & 6 \\
\texttt{learning\_rate} & 0.05 \\
\texttt{subsample} & 0.8 \\
\texttt{colsample\_bytree} & 0.8 \\
\texttt{min\_child\_weight} & 20 \\
\texttt{reg\_alpha} & 0.1 \\
\texttt{reg\_lambda} & 1.0 \\
\texttt{early\_stopping\_rounds} & 50 \\
\bottomrule
\end{tabular}
\end{table}

\subsection{Play Success Model (Binary Classification)}

\textbf{Target:} $\text{success}_t = \mathbf{1}[\text{EPA}_t > 0] \in \{0,1\}$.\\
\textbf{Loss:} Binary cross-entropy (log-loss).\\
\textbf{Output:} $p(\text{success} \mid x) \in [0,1]$.

The success model complements the EPA regressor. Where EPA gives a signed
point estimate of value, the success model gives a calibrated probability.
This is useful for comparing situations: on 3rd \& 8, what fraction of
plays are successful?

\subsection{Win Probability Model (Binary Classification)}

\textbf{Target:} $\text{win}_t = \mathbf{1}[\text{posteam wins game}] \in \{0,1\}$.\\
\textbf{Loss:} Binary cross-entropy.\\
\textbf{Output:} $p(\text{win} \mid \text{game state}) \in [0,1]$.

The win probability model requires the final game outcome to be joined to
each play-level row. The nflfastR \texttt{result} column (final point
differential) provides this; the target is:
\begin{equation}
  \text{win}_t = \mathbf{1}[\text{result}_t > 0]
\end{equation}
A critical quality metric for WP models is \emph{calibration}: a state
assigned $p = 0.70$ should win approximately 70\% of the time in the
data. Reliability diagrams (calibration curves) are used to validate this.

\subsection{Red Zone TD Model (Binary Classification)}

\textbf{Target:} $\text{td}_t = \mathbf{1}[\text{touchdown on this play}] \in \{0,1\}$.\\
\textbf{Training data:} Filtered to $\text{yardline\_100} \leq 20$ (15,103 plays).\\
\textbf{Loss:} Binary cross-entropy; AUC used as evaluation metric.

Because the red zone is qualitatively different from open-field play
(compressed field, different optimal strategies, higher TD base rate),
a dedicated model trained only on red-zone plays outperforms a general model
applied to the red zone. The RZ dataset is smaller (15k vs 101k plays),
so regularisation is tightened: \texttt{max\_depth=4}, \texttt{min\_child\_weight=15}.

\subsection{Drive Outcome Model (Multi-class Classification)}

\textbf{Target:} Drive result $\in \{\text{touchdown}, \text{field goal}, \text{punt}, \text{turnover}, \text{other}\}$.\\
\textbf{Loss:} Multiclass log-loss (\texttt{multi:softprob}).\\
\textbf{Training data:} First play of each drive (drive-start state).

The drive model predicts the probability distribution over outcomes before
a drive begins. Its features are the drive-start game state (field position,
score, time, timeouts). This enables questions like: ``given that we start
this drive at our own 25-yard line down by 7 with 4 minutes left, what is
the probability distribution over drive outcomes?''

\textbf{Note:} This model requires the \texttt{drive} and
\texttt{fixed\_drive\_result} columns, which are included in the R feature
export from the second pipeline run onward.

% ─────────────────────────────────────────────────────────────────────────────
%  7. TRAINING METHODOLOGY
% ─────────────────────────────────────────────────────────────────────────────
\section{Training Methodology}
\label{sec:training}

\subsection{Why Time-Based Cross-Validation}

A recurring mistake in sports analytics is using random $k$-fold cross-validation
on time-series data. This creates \emph{target leakage}: future observations
appear in the training set and past observations appear in the validation set.
For NFL data this is especially problematic because:
\begin{itemize}
  \item Rule changes (e.g., the 2022 kickoff rule) alter EP values;
  \item Teams and personnel evolve season over season;
  \item A model that has ``seen'' games from 2023 during training will
    report optimistically biased performance estimates when evaluated on 2023 games.
\end{itemize}

The correct evaluation scheme trains on all data up to time $T$ and evaluates
on data after time $T$. Here: train on 2021--2022, validate on 2023. This simulates
the realistic deployment scenario where the model is used to predict plays in
a season it has never seen.

\subsection{Early Stopping}

Rather than fixing the number of trees to a constant, we use XGBoost's
early stopping: training halts when validation loss fails to improve for
a configurable patience window (50 rounds for regression, 30--40 for
classifiers). This prevents overfitting without manual tuning of
\texttt{n\_estimators}.

The \texttt{set\_params(early\_stopping\_rounds=N)} call configures the stopping
criterion before each \texttt{fit()} invocation. The best model (lowest
validation loss) is automatically preserved.

\subsection{Training Orchestration}

The \texttt{python/training/trainer.py} module provides a CLI:
\begin{lstlisting}[language=bash]
python3 -m python.training.trainer --model all
\end{lstlisting}
Each training run:
\begin{enumerate}[noitemsep]
  \item Loads \texttt{play\_features.csv};
  \item Calls \texttt{add\_features(df)} to produce the engineered feature set;
  \item Constructs the time-based split;
  \item Casts feature DataFrames to \texttt{float64} numpy arrays (XGBoost 3.x compatibility);
  \item Fits the model with early stopping;
  \item Computes train and validation metrics;
  \item Saves a JSON report to \texttt{models/reports/<model>.json} including
    training metrics and feature importances;
  \item Serialises the fitted model to \texttt{models/saved/<model>.pkl} via
    \texttt{joblib}.
\end{enumerate}

% ─────────────────────────────────────────────────────────────────────────────
%  8. RESULTS
% ─────────────────────────────────────────────────────────────────────────────
\section{Results and Analysis}
\label{sec:results}

\subsection{EPA Model}

\begin{table}[H]
\centering
\caption{EPA model performance}
\begin{tabular}{lrrrr}
\toprule
Split & RMSE & MAE & R\textsuperscript{2} & Trees \\
\midrule
Train (2021--22) & 1.351 & 0.937 & 0.035 & 122 \\
Val (2023) & 1.384 & 0.955 & 0.010 & --- \\
\bottomrule
\end{tabular}
\end{table}

At first glance, $\text{R}^2 = 0.010$ on the validation set appears poor.
This is correct but \emph{not a failure of the model}. Understanding why
requires thinking carefully about what EPA measures.

EPA is a random variable. The same game state---1st \& 10 at the opponent's
40-yard line---can produce a 40-yard touchdown run ($\text{EPA} \approx +4$),
a routine five-yard gain ($\text{EPA} \approx +0.5$), an incomplete pass
($\text{EPA} \approx -0.5$), or a fumble returned for a touchdown
($\text{EPA} \approx -6$). All of these outcomes are plausible from the
same pre-snap state. The variance in EPA is \emph{mostly outcome variance},
not input variance.

The $\text{R}^2$ metric measures the fraction of total variance explained
by the model. Since the overwhelming majority of EPA variance is driven by
the unpredictable outcome of the play itself (not the game state), the
explainable fraction is small by design. The nflfastR EP model, which was
trained on 20+ years of data, achieves RMSE around $0.80$--$0.90$ on EP
(the level, not the change), which is a fundamentally easier task since EP
is smoother. EPA is the derivative of EP with respect to a discrete
transition, making it inherently noisier.

The model \emph{is} doing something useful: it assigns systematically
higher expected EPA to 1st \& 10 at the opponent's 20 than to 3rd \& 15
at one's own 20. This signal is captured by positive-EPA predictions in
favourable situations and negative-EPA predictions in adverse ones. The
value of the model lies in comparing teams' actual EPA to model-predicted
EPA across many plays (the law of large numbers brings the noise down),
not in predicting any individual play.

\subsection{Play Success Model}

\begin{table}[H]
\centering
\caption{Success model performance}
\begin{tabular}{lrrrr}
\toprule
Split & AUC & Brier Score & Log-Loss & Trees \\
\midrule
Train & 0.642 & 0.231 & 0.654 & 191 \\
Val (2023) & 0.614 & 0.235 & 0.661 & --- \\
\bottomrule
\end{tabular}
\end{table}

AUC $= 0.614$ means the model can distinguish successful from unsuccessful
plays better than random (0.50) in about 61.4\% of pair comparisons. The
Brier score of 0.235 can be compared to the null model (always predicting
the base rate $\bar{p} \approx 0.46$): Brier$_{\text{null}} = 0.46 \times 0.54 \approx 0.248$.
Our model improves on this by about 5\%, consistent with game state providing
modest but real predictive signal.

\subsection{Red Zone TD Model}

\begin{table}[H]
\centering
\caption{Red zone model performance}
\begin{tabular}{lrrr}
\toprule
Split & AUC & Brier Score & Log-Loss \\
\midrule
Train (RZ plays) & 0.782 & 0.131 & 0.415 \\
Val (2023 RZ) & 0.766 & 0.132 & 0.417 \\
\bottomrule
\end{tabular}
\end{table}

The red zone model shows stronger discrimination (AUC 0.766 vs 0.614 for
success). This is because TD probability varies much more steeply with
game state in the red zone: at 1st \& goal from the 1 vs 3rd \& goal
from the 15, the situational features have much greater predictive power.
The extremely small train/val Brier gap (0.131 vs 0.132) suggests the model
generalises well to the 2023 season.

\subsection{Feature Importance}

Based on gain-based feature importance from the EPA model, the most
important predictors are:

\begin{enumerate}
  \item \texttt{yardline\_100} (0.234) --- field position drives expected scoring opportunity
  \item \texttt{down} (0.189) --- 1st down is structurally different from 4th down
  \item \texttt{ydstogo} (0.162) --- yards needed determines the value of a stop
  \item \texttt{game\_seconds\_remaining} (0.108) --- late-game dynamics alter EP nonlinearly
  \item \texttt{score\_differential} (0.097) --- score state changes offensive strategy
  \item \texttt{half\_seconds\_remaining} (0.082) --- two-minute drill effects
  \item \texttt{down\_x\_ydstogo} (0.058) --- the interaction term captures conditional difficulty
\end{enumerate}

These rankings are intuitive: field position and down-and-distance are the primary
determinants of EP. The presence of interaction terms in the top ten confirms that
the engineered features are contributing signal beyond their individual components.

% ─────────────────────────────────────────────────────────────────────────────
%  9. SYSTEM ARCHITECTURE
% ─────────────────────────────────────────────────────────────────────────────
\section{System Architecture}
\label{sec:system}

The pipeline is divided into three layers with clean interfaces between them.

\subsection{R Ingestion and Metric Layer}

The R layer handles all interaction with the nflfastR ecosystem.
\texttt{R/pipeline/run\_pipeline.R} orchestrates the full flow:

\begin{lstlisting}[language=R, caption={run\_pipeline.R structure}]
pbp_raw <- load_pbp(SEASONS)          # nflfastR API call
pbp     <- filter_regular_season(pbp_raw)

# aggregate metrics (team/season level)
save_output(compute_epa(pbp),              "epa_by_team.csv")
save_output(compute_success_rate(pbp),     "success_rate_by_team.csv")
save_output(compute_wpa(pbp),              "wpa_by_team.csv")
save_output(compute_redzone_efficiency(pbp), "redzone_efficiency.csv")
save_output(compute_drive_efficiency(pbp), "drive_efficiency.csv")

# play-level ML dataset
save_output(compute_features(pbp),         "play_features.csv")
\end{lstlisting}

Each \texttt{compute\_*(pbp)} function takes the full play-by-play DataFrame
and returns a summarised DataFrame. This pattern makes each metric independently
testable and extensible.

\subsection{Python Machine Learning Layer}

\begin{verbatim}
python/
  features/
    engineering.py    add_features(), get_feature_cols(), time_split()
    selection.py      top_features(), shap_summary()
  models/
    base.py           NFLModel: fit/predict/evaluate/save/load + _np()
    epa_model.py      XGBRegressor, 600 trees
    success_model.py  XGBClassifier, binary
    rz_model.py       XGBClassifier, RZ subset
    wp_model.py       XGBClassifier, with win-outcome join
    drive_model.py    XGBClassifier, multi:softprob
  training/
    trainer.py        CLI orchestration, reports, serialisation
  export/
    json_exporter.py  metrics.json + prediction_grid.json
  api/
    main.py           FastAPI: /metrics/, /predict/play, /models/
  pipeline/
    loader.py         load_metric(), load_all_metrics()
\end{verbatim}

\subsection{Data Flow Diagram}

\begin{verbatim}
nflfastR
    |
    v
R/pipeline/run_pipeline.R
    |
    +---> data/outputs/epa_by_team.csv
    |     data/outputs/success_rate_by_team.csv
    |     data/outputs/wpa_by_team.csv
    |     data/outputs/redzone_efficiency.csv
    |     data/outputs/drive_efficiency.csv
    |
    +---> data/outputs/play_features.csv
                  |
                  v
    python/features/engineering.py
                  |
                  v
    python/training/trainer.py
                  |
          +-------+-------+-------+-------+
          |       |       |       |       |
          v       v       v       v       v
        EPA   Success    RZ      WP    Drive
        .pkl    .pkl    .pkl    .pkl    .pkl
          |       |       |
          +-------+-------+
                  |
                  v
    python/export/json_exporter.py
                  |
          +-------+-------+
          |               |
          v               v
    web/data/           web/data/
    metrics.json    prediction_grid.json
\end{verbatim}

\subsection{FastAPI Backend}

The \texttt{python/api/main.py} module exposes a REST API:
\begin{table}[H]
\centering
\caption{API endpoints}
\begin{tabular}{lll}
\toprule
Method & Path & Description \\
\midrule
GET & \texttt{/health} & Liveness check \\
GET & \texttt{/metrics/\{metric\}} & Aggregate metric data, filterable by team/season \\
GET & \texttt{/teams/\{team\}/profile} & All metrics for a given team \\
POST & \texttt{/predict/play} & EPA + success + RZ TD predictions for a game state \\
GET & \texttt{/models/feature-importance/\{name\}} & Feature importance from training report \\
GET & \texttt{/models/report/\{name\}} & Full training metrics and hyperparameters \\
\bottomrule
\end{tabular}
\end{table}

The \texttt{/predict/play} endpoint accepts a JSON body:
\begin{lstlisting}[language=Python, caption={Play prediction request schema}]
class PlayState(BaseModel):
    down: int                       # 1-4
    ydstogo: int                    # 1-99
    yardline_100: int               # 1-99
    score_differential: int         # negative = trailing
    half_seconds_remaining: int     # 0-1800
    game_seconds_remaining: int     # 0-3600
    posteam_timeouts_remaining: int # 0-3
    defteam_timeouts_remaining: int # 0-3
    play_type: str                  # "pass" | "run"
\end{lstlisting}

% ─────────────────────────────────────────────────────────────────────────────
%  10. WEB APPLICATION
% ─────────────────────────────────────────────────────────────────────────────
\section{Web Application}
\label{sec:webapp}

\subsection{Design Constraints}

The web page (\texttt{web/nfl.html}) is a \emph{static} single-file HTML
document hosted on GitHub Pages. This imposes a hard constraint: no
server-side Python. Predictions must come from pre-computed data or from
JavaScript run in the browser.

\subsection{Data Loading Strategy}

On page load, the JavaScript attempts to fetch \texttt{./data/metrics.json}.
If this request fails (404 or network error), the page falls back to a
comprehensive inline sample dataset that covers all 32 NFL teams across
three seasons, derived from the same formulas as the real model output.
This means the page is always functional regardless of whether the pipeline
has been run.

When \texttt{data/metrics.json} is present, it overrides the sample data
and the page header indicates the generation timestamp.

\subsection{Play Predictor}

The play predictor section demonstrates real-time predictions without a
running backend:
\begin{enumerate}
  \item \textbf{With \texttt{prediction\_grid.json}:} The exporter pre-computes
    XGBoost predictions on a grid of 5,200 game states
    (4 downs $\times$ 14 distances $\times$ 20 yard lines $\times$ 5 score
    differentials $\times$ median time/timeout values). For any user-entered
    state, the JavaScript finds the nearest grid point (by Manhattan distance
    in normalised coordinates) and returns the pre-computed EPA and success
    probability. This gives exact model predictions with zero backend infrastructure.
  \item \textbf{Without the grid:} A built-in JavaScript approximation using
    the canonical EP table (Burke, 2009) and empirically-derived down/distance
    adjustments provides reasonable estimates.
\end{enumerate}

\subsection{Sections}

\begin{description}[leftmargin=2em, itemsep=0pt]
  \item[League Rankings] Sortable table of all 32 teams by any of six metrics
    (EPA/play, success rate, WPA, RZ TD\%, drive scoring rate, composite rank).
    Horizontal bar chart shows EPA/play for the selected season.
  \item[Team Profile] Four bar charts showing a selected team's EPA, success
    rate, red zone performance, and drive efficiency across the three seasons.
  \item[EPA Model] Model performance cards for all five models, feature
    importance bar chart, EPA alpha chart (each team's over/under-performance
    vs. the model's expectation for their game states), and an EPA vs. success
    rate scatter showing team-level correlation.
  \item[Play Predictor] Form accepting the six primary game-state inputs.
    Returns predicted EPA, success probability, and (in red zone) TD probability.
\end{description}

% ─────────────────────────────────────────────────────────────────────────────
%  11. DISCUSSION
% ─────────────────────────────────────────────────────────────────────────────
\section{Discussion}
\label{sec:discussion}

\subsection{On the EPA Model's R\textsuperscript{2}}

It is worth repeating: a low $\text{R}^2$ at the play level does not mean the
EPA model is uninformative. There is a useful analogy to financial returns
models. A factor model for daily stock returns might explain $\text{R}^2 = 0.03$
of individual stock daily returns, yet be extremely useful for portfolio
construction because it captures systematic risk premia. Similarly, our EPA
model captures the systematic relationship between game state and expected scoring
value. The residuals, averaged over hundreds of plays per team per season, provide
a robust estimate of offensive efficiency above expectation.

\subsection{On Three Seasons of Data}

Training on three seasons (approximately 100,000 plays) is modest by machine
learning standards. The early stopping at 122 trees (rather than the 600
configured) and the mild training/validation gap (RMSE 1.351 vs 1.384) suggest
the model is well-regularised and not significantly overfitting. Adding
historical seasons back to 2010 or further would improve model stability but
introduces complications: rule changes (the 2015 PAT rule change, the 2022
kickoff rule) make older data less representative of current football.

\subsection{On Interpretability}

XGBoost's built-in feature importance (gain-based) measures how much each
feature reduces the objective when used in a split, averaged across all trees.
A more principled approach is SHAP (SHapley Additive exPlanations), which
distributes each prediction additively among the input features in a game-theoretic
framework. The training notebooks include SHAP analysis; the \texttt{shap}
Python library is integrated into the training workflow and supports both
global (summary plot) and local (waterfall plot per prediction) explanations.

% ─────────────────────────────────────────────────────────────────────────────
%  12. FUTURE WORK
% ─────────────────────────────────────────────────────────────────────────────
\section{Future Work}
\label{sec:future}

\subsection{Complete Drive and Win Probability Models}

The drive outcome and win probability models require additional columns
(\texttt{drive}, \texttt{fixed\_drive\_result}, \texttt{result}) that are now
included in the R feature export. Re-running \texttt{run\_pipeline.R} and
retraining will complete these two models.

\subsection{Player-Level Metrics via Model Residuals}

With a trained EPA model, play-level residuals can be attributed to players.
If quarterback $q$ is responsible for plays $\mathcal{P}_q$, their
\emph{quarterback value above expectation} is:
\begin{equation}
  \text{QVAE}(q) = \frac{1}{|\mathcal{P}_q|} \sum_{t \in \mathcal{P}_q}
    (\text{EPA}_t - \hat{\text{EPA}}(x_t))
\end{equation}
This removes situational bias: a quarterback who consistently operates behind
the sticks (3rd \& long) but still produces positive residuals is more impressive
than one who only faces favourable situations.

\subsection{Hyperparameter Optimisation}

The current hyperparameters were set manually based on common defaults for
gradient boosting on tabular data. A Bayesian optimisation pass using
\texttt{optuna} over the search space $\{$\texttt{max\_depth}$\in[3,8]$,
\texttt{learning\_rate}$\in[0.01,0.15]$,
\texttt{subsample}$\in[0.6,1.0]$,
\texttt{colsample\_bytree}$\in[0.5,1.0]\}$ would likely reduce validation
RMSE by 2--5\%.

\subsection{Real-Time Prediction}

The FastAPI backend is already structured for real-time prediction.
The next step is adding a WebSocket endpoint that consumes a live
play-by-play stream and pushes predictions to connected clients. Potential
data sources include NFL's official Next Gen Stats API or scraped ESPN
game-log data.

\subsection{Expanded Season Coverage}

Extending the training window to 2015--2023 (approximately 400,000 plays)
would provide substantially more signal for rare game states (4th-quarter
desperation plays, extreme score differentials). Season indicators or a
rolling time window would handle rule change discontinuities.

% ─────────────────────────────────────────────────────────────────────────────
%  REFERENCES
% ─────────────────────────────────────────────────────────────────────────────
\section*{References}
\addcontentsline{toc}{section}{References}

\begin{enumerate}[label={[\arabic*]}, leftmargin=2.5em, itemsep=2pt]

  \item Carter, V. and Machol, R. E. (1971). ``Operations research on football.''
    \textit{Operations Research}, 19(2), 541--545.

  \item Romer, D. (2006). ``Do firms maximize? Evidence from professional football.''
    \textit{Journal of Political Economy}, 114(2), 340--365.

  \item Burke, B. (2009). ``Expected points.''
    Advanced NFL Stats. \url{http://www.advancedfootballanalytics.com}

  \item Chen, T. and Guestrin, C. (2016). ``XGBoost: A scalable tree boosting system.''
    \textit{Proceedings of the 22nd ACM SIGKDD International Conference on
    Knowledge Discovery and Data Mining}, 785--794.

  \item Friedman, J. H. (2001). ``Greedy function approximation: A gradient
    boosting machine.'' \textit{Annals of Statistics}, 29(5), 1189--1232.

  \item Ho, T. K. (1995). ``Random decision forests.''
    \textit{Proceedings of the 3rd International Conference on Document
    Analysis and Recognition}, 278--282.

  \item Carl, S. and Baldwin, B. (2020). ``nflfastR: Functions to efficiently
    scrape NFL play-by-play data.'' R package.
    \url{https://www.nflfastr.com}

  \item Lundberg, S. M. and Lee, S.-I. (2017). ``A unified approach to
    interpreting model predictions.'' \textit{Advances in Neural Information
    Processing Systems}, 30.

  \item Breiman, L. (2001). ``Random forests.''
    \textit{Machine Learning}, 45(1), 5--32.

  \item Lock, D. and Nettleton, D. (2014). ``Using random forests to estimate
    win probability before each play of an NFL game.''
    \textit{Journal of Quantitative Analysis in Sports}, 10(2), 197--205.

\end{enumerate}

\end{document}
